% 主板
% 主板.tex

\documentclass[12pt,UTF8]{ctexbook}

% 设置纸张信息。
\input{../../../Book/common/page}

% 设置字体，并解决显示难检字问题。
\xeCJKsetup{AutoFallBack=true}
\setCJKfamilyfont{hei}{SimHei}
\setCJKfamilyfont{kai}{KaiTi}

% 目录 chapter 级别加点（.）。
\usepackage{titletoc}
\titlecontents{chapter}[0pt]{\vspace{3mm}\bf\addvspace{2pt}\filright}{\contentspush{\thecontentslabel\hspace{0.8em}}}{}{\titlerule*[8pt]{.}\contentspage}

% 支持目录点击跳转
\usepackage[colorlinks,linkcolor=blue,citecolor=blue,urlcolor=blue]{hyperref}

% 设置 part 和 chapter 标题格式。
\ctexset{
	part/name= {第,卷},
	part/number={\chinese{part}},
	chapter/name={第,篇},
	chapter/number={\arabic{chapter}}
}

% 图片相关设置。
\usepackage{graphicx}
\graphicspath{{Images/}}

% 设置署名格式。
\newenvironment{shuming}{\hfill\zihao{4}}

% 注脚每页重新编号，避免编号过大。
\usepackage[perpage]{footmisc}

% 设置段落间距
\setlength{\parskip}{0.8em plus 0.2em minus 0.1em}

% 列表格式支持，列表项向右偏移。
\usepackage{enumitem}

\usepackage{tabularx}
\usepackage{float}

\title{\heiti\zihao{0} 主板}
\author{WangFei}
\date{\today}

\begin{document}

\maketitle
\tableofcontents

\frontmatter
\chapter{前言}

本书专为完全没有计算机硬件基础的学习者设计，将带你从零开始了解主板的基本概念、组成结构和实际应用。

主板是计算机的"骨架"，所有硬件都通过它连接在一起。本书将涵盖以下内容：
\begin{itemize}
    \item 主板的基本概念和作用
    \item 主板的组成结构和关键元件
    \item 主板的类型和规格
    \item 主板的选购和安装
    \item BIOS/UEFI基础设置
\end{itemize}

\mainmatter

\part{主板基础概念}

\chapter{什么是主板}

主板（Motherboard）是计算机中最大的电路板，它连接了CPU、内存、硬盘、显卡等所有硬件组件，是计算机的"中枢神经系统"。

\section{主板的作用}

主板主要负责以下任务：

\begin{itemize}[leftmargin=2em]
    \item \textbf{硬件连接}：提供CPU、内存、显卡等硬件的插槽和接口
    \item \textbf{数据传输}：通过总线连接各个硬件，实现数据交换
    \item \textbf{供电管理}：分配电源给各个硬件组件
    \item \textbf{控制协调}：通过BIOS/UEFI控制硬件工作
\end{itemize}

\section{主板的外观}

主板通常是一块长方形的电路板，上面布满了各种元件：

\begin{enumerate}
    \item \textbf{CPU插槽}：安装CPU的位置
    \item \textbf{内存插槽}：安装内存模块
    \item \textbf{PCIe插槽}：安装显卡、扩展卡
    \item \textbf{SATA接口}：连接硬盘、光驱
    \item \textbf{芯片组}：主板的控制核心
    \item \textbf{供电电路}：为CPU和内存提供稳定电源
    \item \textbf{接口}：USB、音频、网络等外部接口
\end{enumerate}

\section{主板与其他硬件的关系}

\begin{table}[H]
    \centering
    \begin{tabularx}{\textwidth}{|p{3cm}|X|}
        \hline
        \multicolumn{1}{|c|}{\textbf{硬件}} & \multicolumn{1}{c|}{\textbf{与主板的关系}} \\
        \hline
        CPU & 通过CPU插槽连接，由主板供电和控制 \\
        \hline
        内存 & 通过内存插槽连接，与CPU通过内存总线通信 \\
        \hline
        显卡 & 通过PCIe插槽连接，获取高速带宽 \\
        \hline
        硬盘 & 通过SATA或M.2接口连接，存储数据 \\
        \hline
        电源 & 通过24针ATX接口为主板供电 \\
        \hline
        机箱 & 通过铜柱固定主板，前置面板通过排线连接 \\
        \hline
    \end{tabularx}
    \caption{主板与其他硬件的关系}
    \label{tab:motherboard_relations}
\end{table}

\chapter{主板的历史演进}

\section{早期主板}

\begin{enumerate}
    \item \textbf{ISA总线时代（1980年代）}：
    \begin{itemize}[leftmargin=2em]
        \item 早期扩展总线，速度较慢（8MHz）
        \item 用于连接声卡、网卡等设备
    \end{itemize}
    
    \item \textbf{PCI总线时代（1990年代）}：
    \begin{itemize}[leftmargin=2em]
        \item 32位总线，速度提高到33MHz
        \item 取代ISA成为主流扩展接口
    \end{itemize}
\end{enumerate}

\section{现代主板}

\begin{enumerate}
    \item \textbf{PCIe时代（2000年代至今）}：
    \begin{itemize}[leftmargin=2em]
        \item 串行总线，速度远快于PCI
        \item PCIe 1.0/2.0/3.0/4.0/5.0逐步升级
        \item 主要用于显卡和高速扩展卡
    \end{itemize}
    
    \item \textbf{一体化趋势}：
    \begin{itemize}[leftmargin=2em]
        \item 集成声卡、网卡、显卡
        \item M.2接口普及，支持高速SSD
    \end{itemize}
\end{enumerate}

\part{主板组成结构}

\chapter{核心元件}

\section{CPU插槽}

\begin{enumerate}
    \item \textbf{类型}：
    \begin{itemize}[leftmargin=2em]
        \item Intel：LGA（Land Grid Array）
        \item AMD：PGA（Pin Grid Array）/LGA
    \end{itemize}
    
    \item \textbf{兼容性}：插槽必须与CPU型号匹配
    \item \textbf{安装}：CPU有防呆设计，注意方向
\end{enumerate}

\section{芯片组}

芯片组是主板的"大脑"，负责协调各个硬件：

\begin{enumerate}
    \item \textbf{北桥（已集成）}：
    \begin{itemize}[leftmargin=2em]
        \item 连接CPU、内存、PCIe显卡
        \item 现在已集成到CPU中
    \end{itemize}
    
    \item \textbf{南桥}：
    \begin{itemize}[leftmargin=2em]
        \item 连接硬盘、USB、网卡等低速设备
        \item 现在称为PCH（Platform Controller Hub）
    \end{itemize}
\end{enumerate}

\section{内存插槽}

\begin{enumerate}
    \item \textbf{类型}：DDR4、DDR5（不兼容）
    \item \textbf{数量}：通常2-4个插槽
    \item \textbf{双通道}：插在相同颜色插槽实现双通道
\end{enumerate}

\chapter{扩展插槽与接口}

\section{PCIe插槽}

\begin{table}[H]
    \centering
    \begin{tabular}{|p{2cm}|p{3cm}|p{4cm}|p{3cm}|}
        \hline
        \textbf{版本} & \textbf{单通道带宽} & \textbf{常见用途} & \textbf{插槽长度} \\
        \hline
        PCIe 1.0 & 250 MB/s & 早期设备 & x1/x4/x8/x16 \\
        \hline
        PCIe 2.0 & 500 MB/s & 网卡、声卡 & x1/x4/x8/x16 \\
        \hline
        PCIe 3.0 & 1 GB/s & 显卡、SSD & x1/x4/x8/x16 \\
        \hline
        PCIe 4.0 & 2 GB/s & 高端显卡、SSD & x1/x4/x8/x16 \\
        \hline
        PCIe 5.0 & 4 GB/s & 最新设备 & x1/x4/x8/x16 \\
        \hline
    \end{tabular}
    \caption{PCIe版本对比}
    \label{tab:pcie_versions}
\end{table}

\section{SATA接口}

\begin{enumerate}
    \item \textbf{版本}：SATA 2（3 Gbps）、SATA 3（6 Gbps）
    \item \textbf{用途}：连接机械硬盘、SSD、光驱
    \item \textbf{接口数量}：通常4-6个
\end{enumerate}

\section{M.2接口}

\begin{enumerate}
    \item \textbf{特点}：体积小、速度快
    \item \textbf{协议}：支持SATA和PCIe协议
    \item \textbf{规格}：2280（最常见）、2260、2242等
\end{enumerate}

\section{外部接口}

\begin{enumerate}
    \item \textbf{USB接口}：
    \begin{itemize}[leftmargin=2em]
        \item USB 2.0：速度较慢，用于鼠标、键盘
        \item USB 3.0/3.1：高速接口，用于移动硬盘
        \item USB-C：新型接口，支持正反插
    \end{itemize}
    
    \item \textbf{视频接口}：
    \begin{itemize}[leftmargin=2em]
        \item HDMI：高清视频输出
        \item DisplayPort：更高带宽
        \item VGA：老旧接口，已逐渐淘汰
    \end{itemize}
    
    \item \textbf{音频接口}：
    \begin{itemize}[leftmargin=2em]
        \item 3.5mm音频插孔
        \item 光纤音频接口（SPDIF）
    \end{itemize}
    
    \item \textbf{网络接口}：
    \begin{itemize}[leftmargin=2em]
        \item RJ45：千兆以太网接口
        \item 部分主板集成WiFi和蓝牙
    \end{itemize}
\end{enumerate}

\part{主板类型与规格}

\chapter{主板规格}

\section{ATX规格}

\begin{enumerate}
    \item \textbf{尺寸}：305mm × 244mm
    \item \textbf{特点}：标准尺寸，扩展性强
    \item \textbf{适用}：台式机、工作站
\end{enumerate}

\section{Micro-ATX规格}

\begin{enumerate}
    \item \textbf{尺寸}：244mm × 244mm
    \item \textbf{特点}：比ATX小，扩展性适中
    \item \textbf{适用}：紧凑型机箱、家庭电脑
\end{enumerate}

\section{Mini-ITX规格}

\begin{enumerate}
    \item \textbf{尺寸}：170mm × 170mm
    \item \textbf{特点}：最小规格，扩展性有限
    \item \textbf{适用}：HTPC、迷你电脑
\end{enumerate}

\begin{table}[H]
    \centering
    \begin{tabularx}{\textwidth}{|p{3cm}|X|X|X|}
        \hline
        \multicolumn{1}{|c|}{\textbf{规格}} & \multicolumn{1}{c|}{\textbf{ATX}} & \multicolumn{1}{c|}{\textbf{Micro-ATX}} & \multicolumn{1}{c|}{\textbf{Mini-ITX}} \\
        \hline
        尺寸 & 305×244mm & 244×244mm & 170×170mm \\
        \hline
        PCIe插槽 & 3-4个 & 2-3个 & 1个 \\
        \hline
        内存插槽 & 4个 & 2-4个 & 2个 \\
        \hline
        SATA接口 & 6个 & 4个 & 2-4个 \\
        \hline
        扩展性 & 最强 & 适中 & 有限 \\
        \hline
        适用机箱 & 中塔/全塔 & 中塔/小塔 & ITX机箱 \\
        \hline
    \end{tabularx}
    \caption{主板规格对比}
    \label{tab:motherboard_specs}
\end{table}

\chapter{主板芯片组}

\section{Intel芯片组}

\begin{table}[H]
    \centering
    \begin{tabular}{|p{2cm}|p{3cm}|p{6cm}|}
        \hline
        \textbf{芯片组} & \textbf{支持CPU} & \textbf{主要特点} \\
        \hline
        B系列 & 主流处理器 & 性价比高，功能齐全 \\
        \hline
        H系列 & 入门级处理器 & 简化功能，价格便宜 \\
        \hline
        Z系列 & 高端处理器 & 支持超频，扩展性强 \\
        \hline
        X系列 & 旗舰处理器 & 多CPU支持，极致扩展 \\
        \hline
    \end{tabular}
    \caption{Intel芯片组分类}
    \label{tab:intel_chipsets}
\end{table}

\section{AMD芯片组}

\begin{table}[H]
    \centering
    \begin{tabular}{|p{2cm}|p{3cm}|p{6cm}|}
        \hline
        \textbf{芯片组} & \textbf{支持CPU} & \textbf{主要特点} \\
        \hline
        A系列 & AP处理器 & 集成显卡，入门级 \\
        \hline
        B系列 & Ryzen处理器 & 主流选择，功能均衡 \\
        \hline
        X系列 & Ryzen处理器 & 支持超频，高性能 \\
        \hline
        TRX系列 & Threadripper & 多通道内存，专业工作站 \\
        \hline
    \end{tabular}
    \caption{AMD芯片组分类}
    \label{tab:amd_chipsets}
\end{table}

\part{选购与安装}

\chapter{主板选购指南}

\section{兼容性检查}

\begin{enumerate}
    \item \textbf{CPU兼容性}：确认主板支持的CPU插槽和型号
    \item \textbf{内存兼容性}：支持的内存类型（DDR4/DDR5）和最大容量
    \item \textbf{机箱兼容性}：主板规格必须与机箱匹配
    \item \textbf{扩展需求}：PCIe插槽数量、M.2接口数量
\end{enumerate}

\section{品牌推荐}

\begin{enumerate}
    \item \textbf{华硕（ASUS）}：品质可靠，BIOS功能强大
    \item \textbf{微星（MSI）}：性价比高，游戏主板出色
    \item \textbf{技嘉（Gigabyte）}：稳定耐用，产品线丰富
    \item \textbf{华擎（ASRock）}：性价比之王，适合预算有限
\end{enumerate}

\section{选购建议}

\begin{enumerate}
    \item \textbf{日常办公}：B系列芯片组即可
    \item \textbf{游戏玩家}：Z/X系列芯片组，支持超频
    \item \textbf{内容创作}：多内存插槽，高速存储接口
    \item \textbf{迷你电脑}：Mini-ITX规格
\end{enumerate}

\chapter{多显卡配置指南}

\section{核心条件总览}

要在主板上成功运行多块显卡（双卡及以上），必须同时满足以下5个层面的要求，缺一不可。

\begin{enumerate}
    \item \textbf{主板：决定"能不能装"和"能跑多快"}
    \begin{itemize}[leftmargin=2em]
        \item \textbf{芯片组与CPU支持}：
        \begin{itemize}[leftmargin=2em]
            \item AMD平台：需X570、X670、X870E等高端芯片组（B系列一般不行）
            \item Intel平台：需Z690、Z790等Z系列芯片组（B、H系列通常不支持拆分）
        \end{itemize}
        \item \textbf{PCIe插槽}：至少要有两条全长PCIe x16物理插槽，且最好是CPU直连
        \item \textbf{插槽间距}：建议$\ge$3槽（约60mm），否则散热困难
        \item \textbf{PCIe通道拆分}：双卡时需支持x8/x8模式，四卡则需x4/x4/x4/x4（只有HEDT平台能实现）
        \item \textbf{BIOS设置}：需开启"Multi-GPU Support"、"Above 4G Decoding"、"Resizable BAR"等选项
        \item \textbf{供电设计}：高端主板会提供双8Pin甚至12VHPWR辅助供电接口
        \end{itemize}
    
    \item \textbf{CPU：提供足够的PCIe直连通道}
    \begin{itemize}[leftmargin=2em]
        \item \textbf{消费级CPU}（Intel Core / AMD Ryzen）：通常只有20~28条直连通道，最多支持2张显卡（拆分后x8/x8），且会占用M.2固态的通道
        \item \textbf{工作站/服务器级CPU}（Intel Xeon / AMD Threadripper）：拥有64~128条直连通道，可以支持4张、7张甚至更多显卡满速运行
    \end{itemize}
    
    \item \textbf{电源：功率和接口必须够用}
    \begin{itemize}[leftmargin=2em]
        \item \textbf{总功率}：需根据所有配件的峰值功耗相加，并预留30\%余量。例如两张RTX 3090建议1200W以上
        \item \textbf{供电线缆}：必须有足够多的PCIe供电线缆（6+2Pin或12VHPWR）连接每张显卡
        \item \textbf{品质}：强烈建议选用80Plus金牌或铂金认证的高品质电源
    \end{itemize}
    
    \item \textbf{散热与机箱：解决热量堆积}
    \begin{itemize}[leftmargin=2em]
        \item \textbf{机箱}：需宽敞、风道优秀，确保冷风能顺畅流经每张显卡
        \item \textbf{显卡散热器类型}：
        \begin{itemize}[leftmargin=2em]
            \item 涡轮风扇：将热气直接排出机箱，适合多卡紧密排列
            \item 开放式散热器：效率高但热风排入机箱，对机箱排风能力要求极高
        \end{itemize}
        \item \textbf{额外散热}：可额外加装风扇强制给显卡缝隙送风
    \end{itemize}
    
    \item \textbf{软件与驱动：最后激活多卡能力}
    \begin{itemize}[leftmargin=2em]
        \item \textbf{系统与驱动}：操作系统和显卡驱动必须为最新版本
        \item \textbf{NVIDIA SLI}：已停止消费级显卡（RTX 30系及以后）的SLI支持，多卡现在主要用于AI计算、渲染，需通过CUDA等接口调用
        \item \textbf{AMD CrossFire}：也已停止RX 7000系及之后的CrossFire支持
        \item \textbf{注意}：如果应用场景（如游戏）不支持多卡并行，买一张更强的单卡远比组双卡划算
    \end{itemize}
\end{enumerate}

\section{深入理解CPU直连及其重要性}

这是多卡搭建中最容易被忽视的瓶颈。

\begin{enumerate}
    \item \textbf{什么是CPU直连？}
    \begin{itemize}[leftmargin=2em]
        \item PCIe插槽直接与CPU内部的控制器通信，不经过主板芯片组（PCH）中转
    \end{itemize}
    
    \item \textbf{为什么必须直连？}
    \begin{itemize}[leftmargin=2em]
        \item 芯片组与CPU之间的上行链路带宽通常只有PCIe 4.0 x4（约8GB/s），而一张显卡的满速带宽需要PCIe 4.0 x16（约32GB/s）
        \item 如果插在走芯片组的插槽上，显卡性能会被严重限制
    \end{itemize}
    
    \item \textbf{如何识别直连插槽？}
    \begin{itemize}[leftmargin=2em]
        \item 查阅主板官网规格，明确写着"CPU直连"的就是
        \item 物理上，最靠近CPU的那根PCIe x16插槽一定是直连的，其余插槽是否直连要看CPU通道总数
    \end{itemize}
    
    \item \textbf{消费级平台的困局}
    \begin{itemize}[leftmargin=2em]
        \item 主流CPU只有20条直连通道（16条给显卡+4条给M.2），所以双卡时必须拆分为x8+x8
        \item 第三张卡只能走芯片组，毫无性能可言
        \item 只有HEDT（工作站级）平台才有足够通道支持3卡以上满速运行
    \end{itemize}
\end{enumerate}

\section{主板具体推荐}

\subsection{消费级双卡平台（Intel LGA1700）}

适合12~14代酷睿处理器：

\begin{table}[H]
    \centering
    \begin{tabular}{|p{4cm}|p{4cm}|p{2cm}|p{2cm}|p{4cm}|}
        \hline
        \textbf{型号} & \textbf{直连插槽规格} & \textbf{供电} & \textbf{参考价（元）} & \textbf{适合场景} \\
        \hline
        华硕ROG MAXIMUS Z790 HERO & 2条PCIe 5.0 x16（x8/x8）+1条PLX拓展 & 18+1相 & 约5499 & 三卡协同/顶级生产力 \\
        \hline
        微星Z790 GAMING PLUS WIFI & 2条PCIe 5.0 x16（x8/x8） & 16+1相 & 约2099 & 稳定双卡性价比首选 \\
        \hline
        华硕Z790-AYW WIFI W & 2条PCIe 5.0 x16（x8/x8） & 14+1相 & 约1899 & 高性价比双卡入门 \\
        \hline
        映泰Z790 VALKYRIE & 2条全速+2条物理x16（走芯片组） & 14+1相 & 约2199 & 需同时插多张加速卡/采集卡 \\
        \hline
    \end{tabular}
    \caption{Intel LGA1700双卡平台主板推荐}
    \label{tab:intel_dual_gpu}
\end{table}

\subsection{消费级双卡平台（AMD AM5）}

适合锐龙7000/9000系列处理器：

\begin{table}[H]
    \centering
    \begin{tabular}{|p{4cm}|p{4cm}|p{2cm}|p{2cm}|p{4cm}|}
        \hline
        \textbf{型号} & \textbf{直连插槽规格} & \textbf{供电} & \textbf{参考价（元）} & \textbf{适合场景} \\
        \hline
        技嘉X870E AORUS MASTER X3D ICE & 3条PCIe 5.0 x16满速 & 18+2相 & 约4799 & 旗舰级三卡满速 \\
        \hline
        微星MPG X870E Carbon Max WiFi & 2条PCIe 5.0 x16（x8/x8） & 16+2相 & 暂无（高端） & 双卡满带宽稳定平台 \\
        \hline
        华硕ProArt B850-Creator WiFi NEO & 2条PCIe 5.0 x16（x8/x8）智能拆分 & 14+2相 & 暂无（专业级） & 内容创作/低延迟协同 \\
        \hline
        华硕TUF GAMING B850M-PLUS WIFI & 3条PCIe x16（带宽分配需查阅） & 14+2相80A & 约1399 & 预算有限但需多卡扩展 \\
        \hline
    \end{tabular}
    \caption{AMD AM5双卡平台主板推荐}
    \label{tab:amd_dual_gpu}
\end{table}

\subsection{工作站/发烧级多卡平台（3卡以上）}

这类平台采用HEDT芯片组，提供海量直连通道，是真正的多卡"赛道"。

\begin{table}[H]
    \centering
    \begin{tabular}{|p{3cm}|p{4cm}|p{3cm}|p{4cm}|}
        \hline
        \textbf{平台} & \textbf{型号} & \textbf{支持显卡数} & \textbf{参考价（元）} \\
        \hline
        AMD Threadripper & 华硕ROG STRIX TRX40-E GAMING & 3张x16满速 & 约4799 \\
        \hline
        AMD Threadripper & 华硕Pro WS WRX80E-SAGE SE WIFI & 4张PCIe 5.0 x16 & 约16999 \\
        \hline
        Intel Xeon & 华硕Pro WS W790-ACE & 3组PCIe 5.0 x16 & 需查询（高端） \\
        \hline
        Intel Xeon & 超微X13SWA-TF & 6条PCIe 5.0 x16（可装4~6卡） & 需查询（极高端） \\
        \hline
    \end{tabular}
    \caption{工作站级多卡平台主板推荐}
    \label{tab:workstation_multi_gpu}
\end{table}

\section{选购时需额外确认的细节}

\begin{enumerate}
    \item \textbf{确认插槽分配}：看主板规格表，确认"CPU直连"插槽数量及拆分模式（如x8/x8）
    \item \textbf{评估显卡厚度}：若显卡为2.5~3槽厚，务必选择插槽间距$\ge$3槽的主板
    \item \textbf{检查BIOS功能}：是否提供"Above 4G Decoding"、"Resizable BAR"、"SR-IOV"等选项（后两者对于虚拟化和直通很重要）
    \item \textbf{供电余量}：如果打算上双4090/5090，建议主板至少有18相以上供电模组
    \item \textbf{冷却设计}：留意主板M.2散热片是否会与第二张显卡冲突，有时需要拆掉
\end{enumerate}

\section{决策树：快速判断你的需求}

\begin{enumerate}
    \item \textbf{只玩玩游戏}：别组多卡，买一张最强单卡（如RTX 5090）更省心
    \item \textbf{做AI训练/3D渲染，但预算有限}：消费级平台+两张中端卡（如4070 Ti Super $\times$2），选Z790或X870E主板
    \item \textbf{做大型科学计算/深度学习多卡并行}：必须上HEDT平台（Threadripper或Xeon），确保每张卡都能跑满带宽
\end{enumerate}

\section{AI训练场景多卡选购避坑指南}

AI训练场景对多卡配置有特殊要求，以下是常见的坑点和避坑建议：

\begin{enumerate}
    \item \textbf{显存一致性陷阱}
    \begin{itemize}[leftmargin=2em]
        \item \textbf{坑点}：混用不同显存容量或不同显存类型的显卡，导致分布式训练失败
        \item \textbf{避坑}：所有显卡必须使用相同容量的显存，且最好是同一型号、同一批次
        \item \textbf{原因}：PyTorch的DDP（分布式数据并行）要求所有进程的GPU显存一致
    \end{itemize}
    
    \item \textbf{NVLink/NVSwitch的重要性}
    \begin{itemize}[leftmargin=2em]
        \item \textbf{坑点}：忽视GPU间通信带宽，导致多卡训练效率低下
        \item \textbf{避坑}：4卡及以上强烈建议选择支持NVLink/NVSwitch的平台
        \item \textbf{对比}：PCIe 4.0 x16带宽约32GB/s，NVLink可达900GB/s以上
        \item \textbf{适用场景}：模型并行、流水线并行等需要频繁GPU间通信的场景
    \end{itemize}
    
    \item \textbf{PCIe通道分配陷阱}
    \begin{itemize}[leftmargin=2em]
        \item \textbf{坑点}：消费级CPU的PCIe通道被M.2硬盘占用，导致显卡只能工作在x4模式
        \item \textbf{避坑}：AI训练平台尽量使用PCIe扩展卡代替M.2硬盘，或选择不占用CPU直连通道的M.2插槽
        \item \textbf{测试}：安装后用`nvidia-smi topo -m`查看GPU间的连接带宽
    \end{itemize}
    
    \item \textbf{散热系统设计缺陷}
    \begin{itemize}[leftmargin=2em]
        \item \textbf{坑点}：开放式散热器的显卡在多卡紧密排列时互相"烤"
        \item \textbf{避坑}：多卡AI训练建议选择涡轮风扇版本的显卡（数据中心级显卡）
        \item \textbf{散热方案}：
        \begin{itemize}[leftmargin=2em]
            \item 方案一：全涡轮显卡+前后通透机箱
            \item 方案二：开放式散热器+强劲风道+显卡支架
            \item 方案三：分体式水冷（成本较高）
        \end{itemize}
        \item \textbf{温度监控}：AI训练时GPU常满载，建议保持温度$\le$85°C
    \end{itemize}
    
    \item \textbf{电源功率计算失误}
    \begin{itemize}[leftmargin=2em]
        \item \textbf{坑点}：按TDP计算功率，实际训练时功耗远超TDP
        \item \textbf{避坑}：AI训练时GPU功耗可能达到TDP的130\%~150\%
        \item \textbf{计算公式}：总功率 = $\sum$（单卡峰值功耗）$\times$ 1.3（余量）+ CPU功耗 + 其他设备功耗
        \item \textbf{示例}：4张RTX 4090（TDP 450W），实际训练功耗可能达600W/张，总功率需$\ge$3000W
    \end{itemize}
    
    \item \textbf{驱动与CUDA版本兼容性}
    \begin{itemize}[leftmargin=2em]
        \item \textbf{坑点}：最新驱动不一定最稳定，不同显卡型号需要不同驱动版本
        \item \textbf{避坑}：
        \begin{itemize}[leftmargin=2em]
            \item 使用NVIDIA Studio驱动而非Game Ready驱动
            \item 确认PyTorch/TensorFlow支持的CUDA版本
            \item 记录稳定的驱动版本，便于批量部署
        \end{itemize}
        \item \textbf{建议}：CUDA 11.8或12.x版本兼容性较好
    \end{itemize}
    
    \item \textbf{消费级显卡的限制}
    \begin{itemize}[leftmargin=2em]
        \item \textbf{坑点}：使用RTX 30/40系消费级显卡进行长时间AI训练
        \item \textbf{限制}：
        \begin{itemize}[leftmargin=2em]
            \item ECC显存支持：消费级显卡不支持ECC显存，数据完整性无法保证
            \item 保修政策：长时间高负载运行可能影响保修
            \item NVLink支持：消费级显卡通常不支持NVLink
        \end{itemize}
        \item \textbf{建议}：预算允许的话，优先选择数据中心级显卡（如A100、H100、L40S）
    \end{itemize}
    
    \item \textbf{内存容量不足}
    \begin{itemize}[leftmargin=2em]
        \item \textbf{坑点}：忽视CPU内存容量，导致数据加载成为瓶颈
        \item \textbf{避坑}：多卡AI训练建议内存容量 $\ge$ 每张显卡显存的2倍
        \item \textbf{示例}：4张24GB显存的显卡，建议CPU内存 $\ge$ 192GB
        \item \textbf{内存类型}：优先选择高带宽内存（如DDR5-4800及以上）
    \end{itemize}
    
    \item \textbf{BIOS设置遗漏}
    \begin{itemize}[leftmargin=2em]
        \item \textbf{坑点}：未开启必要的BIOS选项导致GPU无法正常工作
        \item \textbf{必须开启}：
        \begin{itemize}[leftmargin=2em]
            \item Above 4G Decoding：支持4GB以上地址空间分配
            \item Resizable BAR：提升GPU显存访问效率
            \item SR-IOV：支持虚拟化直通（适用于容器化部署）
            \item PCIe Power Management：设置为Performance模式
        \end{itemize}
    \end{itemize}
    
    \item \textbf{线缆与接口隐患}
    \begin{itemize}[leftmargin=2em]
        \item \textbf{坑点}：使用转接线或劣质PCIe供电线
        \item \textbf{避坑}：
        \begin{itemize}[leftmargin=2em]
            \item 避免使用12VHPWR转接线（易发热）
            \item 优先使用原生供电线缆
            \item 确保每条PCIe供电线只连接一张显卡
        \end{itemize}
    \end{itemize}
\end{enumerate}

\section{AI训练平台推荐配置}

\begin{table}[H]
    \centering
    \begin{tabular}{|p{3cm}|p{4cm}|p{4cm}|p{4cm}|}
        \hline
        \textbf{预算级别} & \textbf{推荐配置} & \textbf{适用场景} & \textbf{预估成本（元）} \\
        \hline
        入门级（双卡） & Intel i7-14700K + Z790 + 2×RTX 4070 Ti Super & 小规模模型训练、推理 & 约25000 \\
        \hline
        进阶级（四卡） & AMD Threadripper PRO + WRX80E + 4×RTX 4090 & 中等规模模型训练 & 约80000 \\
        \hline
        专业级（四卡） & Intel Xeon W + W790 + 4×NVIDIA L40S & 企业级推理、中等训练 & 约150000 \\
        \hline
        旗舰级（八卡） & AMD Threadripper PRO 7995WX + WRX90E + 8×NVIDIA H100 & 大规模模型训练 & 约1500000 \\
        \hline
    \end{tabular}
    \caption{AI训练平台推荐配置}
    \label{tab:ai_training_configs}
\end{table}

\chapter{主板安装}

\section{安装步骤}

\begin{enumerate}
    \item \textbf{准备工作}：
    \begin{itemize}[leftmargin=2em]
        \item 准备机箱、主板、CPU、内存
        \item 准备螺丝刀、铜柱
    \end{itemize}
    
    \item \textbf{安装铜柱}：
    \begin{itemize}[leftmargin=2em]
        \item 根据主板规格在机箱上安装铜柱
        \item 确保铜柱位置正确
    \end{itemize}
    
    \item \textbf{安装CPU和内存}：
    \begin{itemize}[leftmargin=2em]
        \item 在主板上安装CPU（注意防呆缺口）
        \item 安装内存模块（对准插槽）
    \end{itemize}
    
    \item \textbf{固定主板}：
    \begin{itemize}[leftmargin=2em]
        \item 将主板放入机箱
        \item 使用螺丝固定到铜柱上
    \end{itemize}
    
    \item \textbf{连接线缆}：
    \begin{itemize}[leftmargin=2em]
        \item 连接24针ATX电源线
        \item 连接CPU供电线
        \item 连接前置面板线（PWR\textunderscore SW、RESET\textunderscore SW等）
        \item 连接SATA线、机箱风扇线
    \end{itemize}
\end{enumerate}

\section{注意事项}

\begin{itemize}[leftmargin=2em]
    \item 安装前释放静电（触摸金属物体）
    \item 注意CPU方向，不要用力过猛
    \item 线缆整理，保持机箱内部通风
\end{itemize}

\part{BIOS与设置}

\chapter{BIOS基础}

\section{什么是BIOS}

BIOS（Basic Input/Output System）是主板上的固件，负责：

\begin{enumerate}
    \item \textbf{硬件自检}：开机时检测硬件是否正常
    \item \textbf{引导操作系统}：从硬盘启动系统
    \item \textbf{硬件配置}：设置CPU、内存等参数
\end{enumerate}

\section{进入BIOS的方法}

开机时按特定按键进入BIOS：

\begin{table}[H]
    \centering
    \begin{tabular}{|p{3cm}|p{3cm}|p{6cm}|}
        \hline
        \textbf{品牌} & \textbf{按键} & \textbf{备注} \\
        \hline
        华硕 & Del 或 F2 & 开机时反复按 \\
        \hline
        微星 & Del 或 F2 & 开机时反复按 \\
        \hline
        技嘉 & Del 或 F2 & 开机时反复按 \\
        \hline
        联想 & F1 或 F2 & 部分机型按Enter再按F1 \\
        \hline
        戴尔 & F2 或 F12 & F12可选择启动项 \\
        \hline
    \end{tabular}
    \caption{进入BIOS的按键}
    \label{tab:bios_keys}
\end{table}

\section{UEFI与传统BIOS}

\begin{table}[H]
    \centering
    \begin{tabularx}{\textwidth}{|p{4cm}|X|X|}
        \hline
        \multicolumn{1}{|c|}{\textbf{对比项}} & \multicolumn{1}{c|}{\textbf{传统BIOS}} & \multicolumn{1}{c|}{\textbf{UEFI}} \\
        \hline
        界面 & 蓝色文字界面 & 图形化界面，支持鼠标 \\
        \hline
        启动速度 & 较慢 & 较快 \\
        \hline
        硬盘支持 & 最大2TB & 支持2TB以上大硬盘 \\
        \hline
        安全性 & 较低 & 支持Secure Boot \\
        \hline
        更新方式 & 需要特殊工具 & 可在系统内更新 \\
        \hline
    \end{tabularx}
    \caption{UEFI与传统BIOS对比}
    \label{tab:uefi_vs_bios}
\end{table}

\chapter{常见BIOS设置}

\section{启动顺序设置}

\begin{enumerate}
    \item 进入BIOS的"Boot"选项
    \item 设置第一启动项为硬盘
    \item 可选择U盘、光盘等作为临时启动项
\end{enumerate}

\section{内存频率设置}

\begin{enumerate}
    \item 进入"Memory"或"DRAM"选项
    \item 开启XMP（Intel）或EXPO（AMD）
    \item 选择预设的内存频率配置文件
\end{enumerate}

\section{CPU超频设置}

\begin{enumerate}
    \item 仅适用于支持超频的CPU和主板
    \item 进入"Overclocking"选项
    \item 调整CPU倍频、电压等参数
    \item 注意：超频有风险，可能影响稳定性
\end{enumerate}

\part{维护与故障排查}

\chapter{主板维护}

\section{日常维护}

\begin{enumerate}
    \item \textbf{定期清理}：使用压缩空气清理灰尘
    \item \textbf{检查线缆}：确保所有线缆连接牢固
    \item \textbf{更新BIOS}：定期更新以获得更好的兼容性
\end{enumerate}

\section{常见故障排查}

\begin{enumerate}
    \item \textbf{电脑无法开机}：
    \begin{itemize}[leftmargin=2em]
        \item 检查电源是否连接
        \item 检查CPU和内存是否安装正确
        \item 清除CMOS（取下电池等待几分钟）
    \end{itemize}
    
    \item \textbf{开机无显示}：
    \begin{itemize}[leftmargin=2em]
        \item 检查显示器连接线
        \item 检查显卡是否安装到位
        \item 尝试使用主板集成显卡
    \end{itemize}
    
    \item \textbf{USB接口不工作}：
    \begin{itemize}[leftmargin=2em]
        \item 检查BIOS中USB是否启用
        \item 更新主板芯片组驱动
        \item 尝试其他USB接口
    \end{itemize}
\end{enumerate}

\part{附录}

\chapter{常见问题解答}

\section{常见问题}

\begin{enumerate}
    \item \textbf{Q: 如何确认主板与CPU兼容？}
    \\A: 查看主板官方网站的CPU支持列表，确认CPU型号在支持范围内。
    
    \item \textbf{Q: 主板需要电池吗？}
    \\A: 需要，CMOS电池用于保存BIOS设置和系统时间。
    
    \item \textbf{Q: 什么是CMOS？}
    \\A: CMOS是主板上的一块存储芯片，保存BIOS设置信息。
    
    \item \textbf{Q: 如何更新BIOS？}
    \\A: 可通过BIOS内的更新功能或厂商提供的工具更新。
\end{enumerate}

\chapter{参考资料}

\begin{itemize}[leftmargin=2em]
    \item 华硕官方网站：https://www.asus.com
    \item 微星官方网站：https://www.msi.com
    \item 技嘉官方网站：https://www.gigabyte.com
    \item 华擎官方网站：https://www.asrock.com
\end{itemize}

\backmatter

\end{document}